\subsection{Alternating Linear Forms}

\begin{exercise}{9b1}
    Suppose $m$ is a positive integer.
    Show that $\dim V\mlnset{m} = (\dim V)^m$.
\end{exercise}

\begin{exercise}{9b2}
    Suppose $n \geq 3$ and $\alpha \colon \fbb^n \times \fbb^n \times \fbb^n \to \fbb$ is defined by
    \[
        \alpha((x_1, \ldots, x_n), (y_1, \ldots, y_n), (z_1, \ldots, z_n))
        = x_1 y_2 z_3 - x_2 y_1 z_3 - x_3 y_2 z_1 - x_1 y_3 z_2 + x_3 y_1 z_2 + x_2 y_3 z_1.
    \]
    Show that $\alpha$ is an alternating $3$-linear form on $\fbb^n$.
\end{exercise}

\begin{exercise}{9b3}
    Suppose $m$ is a positive integer and $\alpha$ is an $m$-linear form on $V$ such that $\alpha(v_1, \ldots, v_m) = 0$ whenever $v_1, \ldots, v_m$ is a list of vectors in $V$ with $v_j = v_{j+1}$ for some $j \in \set{1, \ldots, m - 1}$.
    Prove that $\alpha$ is an alternating $m$-linear form on $V$.
\end{exercise}

\begin{exercise}{9b4}
    Prove or give a counterexample: If $\alpha \in V\mlnalt{4}$, then
    \[
        \set{(v_1, v_2, v_3, v_4) \in V^4 : \alpha(v_1, v_2, v_3, v_4) = 0}
    \]
    is a subspace of $V^4$.
\end{exercise}

\begin{exercise}{9b5}
    Suppose $m$ is a positive integer and $\beta$ is an $m$-linear form on $V$.
    Define an $m$-linear form $\alpha$ on $V$ by
    \[
        \alpha(v_1, \ldots, v_m) = \sum_{(j_1, \ldots, j_m) \in \perm m} (\sign(j_1, \ldots, j_m)) \beta(v_{j_1}, \ldots, v_{j_m})
    \]
    for $v_1, \ldots, v_m \in V$.
    Explain why $\alpha \in V\mlnalt{m}$.
\end{exercise}

\begin{exercise}{9b6}
    Suppose $m$ is a positive integer and $\beta$ is an $m$-linear form on $V$.
    Define an $m$-linear form $\alpha$ on $V$ by
    \[
        \alpha(v_1, \ldots, v_m) = \sum_{(j_1, \ldots, j_m) \in \perm m} \beta(v_{j_1}, \ldots, v_{j_m})
    \]
    for $v_1, \ldots, v_m \in V$.
    Explain why
    \[
        \alpha(v_{k_1}, \ldots, v_{k_m}) = \alpha(v_1, \ldots, v_m)
    \]
    for all $v_1, \ldots, v_m \in V$ and all $(k_1, \ldots, k_m) \in \perm m$.
\end{exercise}

\begin{exercise}{9b7}
    Give an example of a nonzero alternating $2$-linear form $\alpha$ on $\rbb^3$ and a linearly independent list $v_1, v_2$ in $\rbb^3$ such that $\alpha(v_1, v_2) = 0$.

    \note{This exercise shows that 9.39 can fail if the hypothesis that $n = \dim V$ is deleted.}
\end{exercise}